\documentclass[11pt]{article}
\usepackage[margin=1in]{geometry}
\usepackage{setspace}
\doublespacing
\usepackage{parskip}
\usepackage{times}
\usepackage{titlesec}
\titleformat{\section}
  {\normalfont\Large\bfseries}{\thesection}{1em}{}
\titleformat{\subsection}
  {\normalfont\large\bfseries}{\thesubsection}{1em}{}
\usepackage{fancyhdr}
\pagestyle{fancy}
\fancyhf{}
\rhead{Moira's Journal: Up the Silverthread}
\lhead{Page \thepage}
\cfoot{}

\begin{document}

\section*{Journal Entry 1: The Threshold}
\textit{Day 3 of the Ascent. Stonehaven's last outpost is behind us---the smoke of their hearths just a memory against the morning mist. The Silverthread narrows here, choked with fallen timber and slick stone where the mountains first claim it. Tam moves ahead with that quiet certainty of his. No hesitation. No over-mapping of each step. I follow, notebook open but often closed as I watch him work.}

The river speaks a language Tam understands without translation. Where I see only tangled roots and shifting currents, he reads the water's mood---the way it hesitates before a submerged log, the slight upwelling that promises a deeper channel, the faint tremor on the surface that warns of cold upwelling from a spring unseen. He doesn't pause to consult his compass or read the stars for omens. He dips a finger in the current, feels its temperature and speed, and adjusts our course with a small shift of his pack.

It's not magic. Not exactly. It's attention honed by years of listening to what the world insists on telling those who will hear it. His ex-partner, that Aeler Edgewalker who taught him, called this "reading the river's breath." I call it practical wisdom---the kind that keeps boots dry and packs light when the current turns treacherous.

\subsection*{Tam's Rule No. 1: The First Read}
\textit{"Trust your first read of a crossing, but watch for the second guess. The river will try to trick you into overthinking. Don't let it."}

We made camp late afternoon where the river bends sharp around a granite shoulder. Tam picked the spot---not from any map, but by watching the water pile on the inside curve. \textit{Calmer}, he said. \textit{Dry ground under the ferns}. I'd have missed it.

He moved with purpose, testing each stone before committing weight, stance wide and ready. When a loose scree shifted under his boot, he didn't stumble---he flowed with it, using the momentum to turn and brace against the rock face. No flourish. Just the efficient use of what the mountain offered. The implied difficulty: footing treacherous (DV 4 equivalent), yet his movement stayed fluid, never slipping into desperation. That's the mark of his training. He meets resistance not with force but with alignment, turning potential falls into adjustments.

\subsection*{Tam's Rule No. 2: Preparation Over Planning}
\textit{"Maps lie. The trail tells the truth if you're quiet enough to hear it. Planning is for those who fear the truth; preparation is for those who respect it."}

When I asked why he doesn't plot our course more carefully, he gave that half-smile of his---knowing, a little weary of my academic caution. Then he said the line I've since written into my own notebook: \textit{"Maps lie, Moira. The river tells the truth if you're quiet enough to hear it."}

He carries his knowledge lightly, ready to discard it the moment the river says otherwise.

\subsection*{Bestiary of the Headwaters (First Pass)}
\begin{itemize}
    \item \textbf{Silverthread Minnow}: Small, silver-scaled fish that dart in tight schools just below the surface. Their sudden, coordinated flashes often indicate subtle changes in current or temperature---useful for reading the river's "mood" without entering the water. Harmless, but their absence can warn of upstream disturbances or predators.
    \item \textbf{Granite Lizard}: Stocky, grey-scaled reptiles that bask on sun-warmed stones. They sense vibrations through rock with uncanny precision, freezing perfectly still when footsteps approach on loose scree. Practical takeaway: their stillness is not fear but assessment. If they bolt, you've made too much noise. If they stay, you can move carefully.
\end{itemize}

\section*{Journal Entry 2: The Still Pool}
\textit{Day 4. The river has widened into a series of shallow pools, each separated by low falls that sound like conversation. Tam is quieter than usual. He says the water is "listening." I'm not sure if he means it literally. I'm not sure it matters.}

We reached a stretch where the Silverthread slows almost to stillness. The pools are clear, deep, and cold---the kind of water that takes your breath if you slip. Tam stopped at the edge of the largest pool, crouched, and stared into the depths for a long count of thirty.

"Someone's been here," he said. "Not recently. But they left something."

He pointed to a dark shape on the bottom, half-hidden by silt and a fallen branch. A pack. Old leather, rotting straps, one buckle still gleaming.

\subsection*{Tam's Rule No. 3: Leave What You Find (Unless You Know Why It Was Left)}
\textit{"A dead traveler's pack is not a gift. It is a message. Read the message before you touch the pack."}

He didn't reach for it. Instead, he walked the perimeter of the pool, scanning the banks, the overhanging branches, the direction of the current. Then he sat on a flat stone and asked me what I saw.

I listed: the pack (obvious). The branch that had fallen across it (natural, but maybe placed). The way the silt was banked higher on one side (current shifts when the river rises). The absence of bones (no body nearby).

"Anything else?"

I looked again. The buckle. It was still bright. Not corroded like the leather. Newer. Or made of something that doesn't tarnish.

"Someone put that pack there recently," I said. "The buckle is fresh. The rest is staged."

Tam nodded. "The river didn't take this traveler. Someone wanted us to think it did. We move on. Don't touch the water."

We camped a mile upstream. No one followed.

\subsection*{Bestiary of the Headwaters (Continued)}
\begin{itemize}
    \item \textbf{Pool-Biter}: A freshwater leech that grows to the length of a hand, pale and almost invisible in still water. It detects warmth and movement. Its bite is not venomous but numbs the skin, making it difficult to notice until you're weak from blood loss. Practical takeaway: If a pool has no minnows, check the edges for these. Minnows avoid them.\\
    \textit{(Tam's addendum: "Also, piss upstream. The smell warns them off." I did not ask how he learned this.)}
\end{itemize}

\section*{Journal Entry 3: The Ford and the Fool}
\textit{Day 6. We crossed the Silverthread at a place Tam called the "Wager's Ford." He said the name came from a bet lost by a surveyor who couldn't decide where the trail went. The ford, he claimed, only appears when the river is at a certain height. It was not at that height. We crossed anyway.}

The water was thigh-high and faster than I liked. Tam went first, testing each step with his staff, never lifting his feet but sliding them forward against the current. "Smooth stones," he called back. "Good grip if you don't rush."

I did not rush. I did not fall. But I felt the current testing my balance, the cold seeping through my boots, the weight of my pack trying to pull me sideways. When I reached the far bank, Tam was already building a fire.

"You hesitated," he said. "Twice. Once when you felt the cold, once when you saw the depth. That's fine. Hesitation is not fear. It's calculation. The problem is when you let calculation become refusal."

\subsection*{Tam's Rule No. 4: The Ford}
\textit{"Do not fight the current. Move with it, angled downstream, and let it push you across. Fighting the current is how you drown. Trusting it is how you cross."}

I asked how he knows which fords are safe. He shrugged. "You don't. You learn which ones are dangerous. The rest you try."

He then gave me a list of signs that a ford is \textit{not} dangerous---things I should have noticed myself:
\begin{itemize}
    \item The bottom is visible. If you cannot see the stones, the river is too stirred up (or too deep).
    \item The current smooths out in the middle, not churning. Churning means a submerged rock or a drop-off.
    \item The far bank has a trail leading away. Not a game trail---a human trail. Someone else crossed here recently.
    \item The water temperature is consistent across the width. Cold patches mean springs or underground channels. Avoid those lines.
\end{itemize}

\subsection*{Tam's Rule No. 5: The Cold Patch}
\textit{"If a patch of water is colder than the rest, do not step there. Springs will pull you under. Also, things that like cold live in springs. They are not friendly."}

He did not elaborate. I did not ask.

\section*{Journal Entry 4: The Watcher in the Pines}
\textit{Day 8. We left the river and climbed into the pine forest that skirts the mountain's western shoulder. The air is thinner here. The trees are old, their trunks scarred by fire and lightning. Tam is more alert than I've seen him. Not afraid---watching.}

We found tracks at midday. Not animal tracks. Boots. Old, worn, the tread pattern unfamiliar. They followed the ridgeline, not the game trail we were using. Parallel. Deliberate.

"Someone is pacing us," Tam said. "Not hunting. Just watching."

He raised his voice slightly. "We're not looking for trouble. Passing through to the high meadows."

No answer. The tracks continued.

We camped that night in a small hollow ringed by boulders. Tam set no fire. We ate cold rations in the dark. He laid out a pattern of small stones around the perimeter---something he'd learned from the Aelinnel, he said. "They measure trespass. If someone crosses this line, I'll know."

I asked what he would do if someone crossed.

"Negotiate," he said. "Then run. Then fight. In that order."

No one crossed. The tracks were gone in the morning.

\subsection*{Bestiary of the Headwaters (Final)}
\begin{itemize}
    \item \textbf{Shadow-Stalker (Rumored)}: No confirmed sightings. Tam says he's seen their tracks---three-toed, front-heavy, with a stride that suggests a loping gait. They follow travelers for days, never attacking, only watching. Some say they are guardians of old boundaries. Others say they are scavengers waiting for you to make a mistake. Practical takeaway: If you are being paced, do not run. Running is the mistake they are waiting for. Walk steadily. Make camp in the open. Build a fire. They dislike fire.
    \item \textbf{Pine-Wraith (No, Not That Kind)}: A fungal growth that parasitizes old pines, causing the bark to peel in long spirals. At night, the exposed wood glows faintly green. Not dangerous, but travelers who camp near infected trees report vivid dreams and a tendency to wake facing the wrong direction. Practical takeaway: Avoid camping within twenty paces of a spiraled pine. If you must, tie yourself to a companion. Waking up lost in the dark is how you die.
\end{itemize}

\section*{Journal Entry 5: The Summit}
\textit{Day 10. We reached the high meadow. The Silverthread's source is a spring that bubbles up from a crack in the granite, cold and clear, feeding a small pool ringed by moss. Tam knelt and drank from his cupped hands. Then he poured a small amount onto the stone at the spring's edge.}

"For the mountain," he said. "A toll. The water remembers who drinks without asking."

I asked what he meant. He didn't answer directly. Instead, he pointed at the pool. "Look. No fish. No insects. Nothing lives here. But the water is clean. The cleanest we've had."

"Then why are there no fish?"

"Because something doesn't want them here. Something that lives downstream, maybe. Or something that lives here and leaves at night." He shrugged. "We're not staying. We don't need to know."

He marked the spring on my map---a small symbol, a circle with a dot, the Aeler mark for "water, but ask first."

\subsection*{Tam's Rule No. 6: The Source}
\textit{"A spring that gives water without life is not a gift. It is a loan. Drink, but pour some back. The mountain is watching."}

We camped in the meadow that night. The stars were sharp and cold. Tam pointed out the constellations the guides never mention---the ones the Edgewalkers use to navigate when the old stars are hidden.

"Those," he said, "are the Unwritten. They don't have names. They just point."

Point to what? I asked.

"Home," he said. "If you know how to read them."

I did not ask which way home was. I suspect he would not have told me.

\section*{Closing Notes: What I Carried Down}
The descent took half the time of the ascent. We followed the river, then left it, then found it again at a lower ford. Tam spoke less. I wrote more. The practical lessons he gave me---rules for fords, warnings about cold patches, the meaning of a spiraled pine---have filled the margins of this journal. I will transcribe them properly when I am home.

But the lesson that stays with me is not written. It is the way he moved through the world: always present, never rushed, treating each step as a negotiation rather than a conquest. The mountain did not yield to him. It simply agreed to let him pass.

That is what I came to learn. That is what I will try to teach.

\subsection*{A Final Bestiary Entry (Added at Tam's Request)}
\begin{itemize}
    \item \textbf{The Scholar (Homo thepyrgosianus academicus)}: Common in libraries, rare on trails. Prone to over-noting, under-watching. Carries heavy packs full of things they will not use. Frequently asks questions that have no answer. Practical takeaway: Travel with one if you must. Their notes may save you later. But do not let them navigate. They will get you lost by the third day.
\end{itemize}

\vspace{2em}
\textit{End of the Silverthread Journal.}
\vspace{1em}
\noindent Moira of Stonehaven\\
\textit{Returned to the outpost, boots drying by the fire, ink still wet.}

\section*{Journal Entry 1: The Edge of the Wood}
\textit{Day 1 of the Forest Passage. We've come down from the high meadows into the basin they call the Hollow Wood---not a true hollow, but a bowl of old-growth forest where the light filters green and the ground never fully dries. Tam stopped at the first tree line and unwrapped a length of red thread from his pack. He tied it to a low branch.}

"The Aelaerem do this," he said. "You leave a thread at the edge. It tells the wood you're coming. It tells the Neighbors that you know the custom."

I asked what happens if you don't leave a thread.

"You get turned around. You find your path looping back to where you started. You hear footsteps that aren't yours." He shrugged. "Nothing angry. Just... territorial."

\subsection*{Tam's Rule No. 7: The Red Thread}
\textit{"Leave a thread at the edge of old forest. If you don't have thread, leave a coin. If you don't have a coin, leave a word spoken aloud---a name, a thanks, a promise to do no harm. The wood has ears."}

The trail was clear at first---a game track widened by generations of travelers. But after an hour, the trees began to crowd closer. The light shifted. Shadows seemed to hold their shape even when I turned my head. I caught myself looking over my shoulder twice, three times. Tam didn't look at all. He just walked.

"Do you feel it?" I asked.

"Feel what?"

"The watching."

He gestured at the canopy. "The trees watch everything. They don't have opinions. They just watch. Stop worrying about being seen and start worrying about what you're carrying."

I checked my pack. I had brought dried meat, hardtack, a compass, a notebook, a knife, and a small iron bell (a gift from a Aeler friend, for warding off "unpleasant attentions").

"The bell," Tam said. "Keep it in your pack. Don't ring it unless I tell you."

"Why?"

"Because bells are loud. Loud is rude. The wood prefers quiet. Ring the bell and the wood will think you're starting an argument."

\subsection*{Bestiary of the Hollow Wood (First Pass)}
\begin{itemize}
    \item \textbf{Hedge-Crawler}: A small, many-legged insect that nests in the bark of old oaks. It emerges at dusk, forming slow-moving columns along branches and fallen logs. Harmless, but if you disturb a column, it releases a bitter-smelling oil that stains skin for days. Practical takeaway: Step over, not through. The oil attracts scavengers.
    \item \textbf{Thread-Bird (Fairy Warbler)}: A finch-sized bird with red markings that resemble knotted thread. It follows travelers through the Hollow Wood, perching ahead on branches, sometimes calling in a pattern that mimics human speech. Local tradition says the birds carry messages from the Neighbors. Tam says they're just curious. Practical takeaway: If a thread-bird mimics a warning ("go back," "turn left"), do not trust it. They mimic what they hear, not what they mean.
\end{itemize}

\section*{Journal Entry 2: The Corpse-Lantern}
\textit{Day 2. We made camp at a dry hollow beneath a massive oak. Tam built no fire. "Too deep in the wood," he said. "Fire attracts things that want warmth." We ate cold and slept against the tree's roots, our backs to the trunk.}

I woke at midnight to a pale light bobbing between the trees. Silent. Not a lantern---no flame, no glass. Just a soft, green-white glow that moved without visible source.

Tam was already awake. He put a hand on my arm. "Corpse-lantern," he whispered. "Fungus. Grows on dead wood. Follows the wind. Not dangerous."

"Then why are you whispering?"

"Because the thing that \textit{eats} the corpse-lantern isn't polite."

We watched the light drift past, fifty paces away. Behind it, a shape moved---low, broad, silent. No details. Just a shift in the darkness where the light had passed.

"It's gone," Tam said after a long minute. "Sleep. I'll watch."

I didn't sleep. I wrote this by feel in the dark.

\subsection*{Tam's Rule No. 8: Night in the Wood}
\textit{"Do not light a fire unless you are cold enough to die. The things that fear fire are not the only things in the wood. The things that love fire are worse."}

At dawn, Tam examined the ground where the light had passed. He found a patch of trampled moss and a single large footprint---three toes, clawed, the heel deep and wide.

"Bear?" I asked.

"Maybe. Or something that walks like a bear and eats glow-fungus. I've seen tracks like this before. In the Valewood, near the old Aelinnel ruins. The rangers there call them \textit{Lamp-Eaters}. They follow the corpse-lanterns and sometimes follow travelers who carry light."

I checked the small lantern in my pack---glass, oil, wick. Useless here.

"Do they attack?"

"They eat light. Not people. But if you have light and they can't eat it, they get frustrated." He made a gesture of apology at the empty woods. "We're not worth the trouble. Keep moving."

\section*{Journal Entry 3: The Counting Stones}
\textit{Day 3. We found a circle of standing stones in a clearing. Seven stones, low and mossy, their tops worn smooth. Aelinnel work, Tam said---old, maybe older than the human settlements in the valley.}

"They used these to count," he said. "Not numbers. Debts. Every stone marked a transaction between the Aelinnel and the forest. Wheat for timber. Salt for safe passage."

"Seven stones," I said. "Why not eight?"

"Because the eighth was the forest's share. They didn't carve it. They just left the space." He pointed to a gap in the circle, a flat patch of earth where another stone might have stood. "That's where the Neighbors sat. Still do, maybe. We should leave something."

We left a pinch of salt on each stone. Tam spoke a phrase in a language I didn't recognize---Aelinnel, he said later, the old counting tongue. "It means 'the debt is acknowledged, not yet paid.'"

\subsection*{Tam's Rule No. 9: Standing Stones}
\textit{"Do not sit on a standing stone. Do not lean against it. Do not carve your name into it. If you must touch it, touch the north face. The south face belongs to the dead."}

I asked how he knew which face was north. He pointed to the moss.

"Moss grows on the north face. Except when it doesn't. Then you have a different problem."

\subsection*{Bestiary of the Hollow Wood (Continued)}
\begin{itemize}
    \item \textbf{Counting-Stone Moss}: A pale green lichen that grows exclusively on Aelinnel standing stones. It glows faintly on nights when the Neighbors are said to visit. Not dangerous, but its presence indicates the stone is still "active"---the forest still honors the old debts. Practical takeaway: If you find a standing stone without moss, do not camp nearby. The debt has been forgotten, and something else may have claimed the space.
    \item \textbf{Lamp-Eater (Valewood bear?)}: A large, nocturnal omnivore that feeds on bioluminescent fungi. It has poor eyesight but an acute sense of smell. It follows sources of light out of confusion, not aggression. If one approaches your camp, extinguish all lights and remain still. It will lose interest and wander away.
\end{itemize}

\section*{Journal Entry 4: The Fork in the Path}
\textit{Day 5. The trail split. One branch continued west, toward the lowlands. The other turned north, climbing into older forest, the trees thicker, the understory choked with thorn and bramble. Tam stopped at the fork and examined both paths for a long time.}

"The north path is shorter," he said. "But it passes through a stretch of ground the Aelaerem avoid. They don't say why."

"The west path?"

"Longer. Safer. But we lose two days."

I asked what he recommended.

He sat on a fallen log. "You're the one writing the guide. What do you see?"

I looked. The north path had more light---the canopy opened slightly, letting in sun. But the ground was wetter, the thorns thicker. The west path was drier, but darker---the trees tightened overhead, and the air stilled.

"The north path is trying to look inviting," I said. "That makes me suspicious."

Tam nodded. "Good. West it is."

\subsection*{Tam's Rule No. 10: The Inviting Path}
\textit{"If one fork looks significantly easier than the other, take the harder one. The forest does not offer charity. It offers traps."}

We took the west path. Within an hour, the darkness deepened, and I began to hear footsteps behind us---not mine, not Tam's. He didn't acknowledge them.

"Don't look back," he said. "They're not following you. They're following the path. Just keep walking."

We walked. The footsteps faded. The path opened. I did not look back.

\section*{Journal Entry 5: The Clearing of Silence}
\textit{Day 6. We emerged from the wood into a clearing with a small pond. The water was black, perfectly still, reflecting the sky like a mirror. No birds. No insects. No wind.}

Tam walked to the water's edge and stood for a long time. Then he sat, cross-legged, and gestured for me to do the same.

"The silence," he said, "is not empty. Listen."

I listened. Nothing. Then, after a minute, a sound: water dripping. Not from the pond---from somewhere above, from the trees, from the sky? The drops fell into the pond, each one a clear note, like a bell.

"The Neighbors are here," Tam whispered. "They want to know why we crossed their wood."

"What do we say?"

"The truth."

He spoke again in that old tongue---slow, careful, each word placed like a stone. I caught fragments: "passage," "debt," "no violence," "thank you." Then he fell silent.

The dripping stopped. The silence returned. Then, from the far side of the pond, a voice. Not human. Not animal. Just... a voice. It said one word, in Utaran.

*"Stay."*

Tam did not move. I did not move. The voice did not speak again.

We stayed for one hour. Then we rose, walked to the far side of the clearing, and found a path we had not seen before---wide, dry, leading down toward the lowlands.

"Now we go," Tam said.

I asked what would have happened if we had left earlier.

"The Neighbors would have considered us rude. They might have closed the path behind us."

"And if we had stayed too long?"

"Then we would have become part of \textit{their} clearing. The silence would have kept us."

\subsection*{Tam's Rule No. 11: The Neighbors}
\textit{"If a clearing is silent, wait. Someone is speaking. You just can't hear them yet. When you hear them, do not argue. Do not bargain. Do not threaten. Answer only the question they ask, and answer truthfully. Then wait to be dismissed."}

\subsection*{A Final Bestiary Entry (For the Road)}
\begin{itemize}
    \item \textbf{The Neighbors}: Not a species. Not a faction. A \textit{category} of attention---the forest's awareness of intruders. They manifest as silence, as footsteps that aren't there, as a voice that speaks your language but does not belong to any throat. They do not harm. They only \textit{remember}. Practical takeaway: Do not test them. Do not ignore them. Acknowledge them politely and move on. The wood is old. It will outlast you. It knows this.
\end{itemize}

\vspace{2em}
\textit{End of the Hollow Wood Journal. We reached the lowlands on the eighth day. Tam untied his red thread from the first tree and coiled it back into his pack. "The wood remembers the thread," he said. "Next time, I won't need a new one."}

\noindent Moira of Stonehaven\\
\textit{Dry ground beneath my feet, sunlight on my face, ink almost gone.}

\section*{Journal Entry 1: The Dying Village}
\textit{Day 1 of the Detour. We emerged from the Hollow Wood into a narrow valley. The village---name scratched on a board at the trailhead: \textit{Grey Hollow}---sat on a low rise above a creek. Smoke rose from a few chimneys, but there was no bustle, no children shouting. The fields were overgrown. The few people we saw moved slowly, heads down.}

Tam didn't speak until we reached the center. A well, a mossy statue of a saint (the Everflame, I think, though the features had worn smooth), and a longhouse with a broken weathervane.

"Something's wrong," he said. "Not plague. Not famine. Worse."

\subsection*{Tam's Rule No. 12: The Weight of a Village}
\textit{"A village that has lost its heart does not starve. It fades. The people forget to laugh. The fields stop yielding. The children stop playing. Something was taken that should not have been taken."}

We were invited to stay by an elder---a woman named Ylla, her hands calloused, her eyes tired but still sharp. She gave us bread, salt, thin soup. The portions were small. She apologized twice.

After the meal, Tam excused himself to speak with Ylla alone. I stayed by the longhouse fire, watching the villagers drift past. A man mended a net with movements so slow and deliberate it looked like ritual. A woman carried water from the creek, spilling half, not noticing.

When Tam returned, his face was still.

"Their luck-relic," he said. "A stone. Shaped like a closed eye. Kept in a shrine above the village. They say it brought good harvest, healthy children, safe passage through the woods. A month ago, a stranger came---a relic hunter, by his gear. He offered to buy it. They refused. That night, he stole it."

"And they chased him?"

"Three men followed him into a cave south of the village. The men did not return. The cave..." Tam paused. "Ylla says it is not a cave. It is a \textit{door}. Old. Older than the mountains. The men who entered did not come out. The stranger did not come out. The relic is still inside, somewhere."

"Why didn't they send for help?"

"They did. No one came. The road from the lowlands has been closed by a rockslide. The passes are watched by things that do not answer to taxes."

Tam stood. "I'll go. You stay here. If I'm not back in three days, leave. Head east to the river. Follow it down."

"No," I said. "I'm going with you."

He looked at me for a long time. Then he nodded. "One bell. Keep it in your pocket. Ring it if you get lost."

\subsection*{Bestiary of Grey Hollow (First Pass)}
\begin{itemize}
    \item \textbf{Hollow-Sick (Condition)}: Not a disease---a spiritual draining caused by the absence of a village's keystone relic. Symptoms include lethargy, forgetfulness, a tendency to wander toward the cave entrance without knowing why. Practical takeaway: If a village offers you hospitality but seems listless, check for their luck-stone. If it's missing, do not accept food or water. The hollow-sick may spread through gifts.
\end{itemize}

\section*{Journal Entry 2: The Door in the Hill}
\textit{Day 2. The cave entrance was not a cave. It was a doorway cut into the limestone, its arch carved with symbols I could not read---not Aelinnel, not Aeler, not any script in my field books. Tam touched the stone. "Warm," he said. "It's been waiting."}

We lit a single lantern. The passage sloped down, then leveled, then widened into a hall. The walls were black basalt, polished smooth. No dust. No cobwebs. No sign of decay.

"Someone built this to last," I whispered.

"Someone built this to \textit{forget}," Tam replied.

The hall ended in a stair that descended into darkness. We counted the steps. Eighty. Then a landing. Then another stair. Eighty more. The air grew cool, still, heavy.

\subsection*{Tam's Rule No. 13: The Deeper You Go}
\textit{"In a dungeon that was built to be forgotten, count your steps. When you reach a multiple of eighty, stop. Listen. The builders marked their thresholds with silence."}

We stopped. Listened. Nothing. Then, from below, a low hum. Not a voice---a vibration, felt through the soles of my boots.

Tam put his hand on the wall. "Pressure change. There's a chamber down there. Something is \textit{living} in it."

We descended.

\section*{Journal Entry 3: The Living Dungeon}
\textit{Day 3. The ziggurat's corridors shifted overnight. I am certain of this. A passage that led east now led north. A stair that descended eighty steps now descended eighty-three. Tam wasn't surprised.}

"The builders didn't want visitors," he said. "They re-arrange the path when someone is inside. Keeps you from finding the center."

"How do you find it, then?"

"You don't. You let the dungeon find you. Walk until the geometry stops fighting. Then ask."

We walked. The air grew warmer. The stone changed from basalt to a pale, veined marble. The walls were carved with scenes---processions, offerings, a figure on a throne that never faced the viewer.

\subsection*{Bestiary of the Ziggurat}
\begin{itemize}
    \item \textbf{Memory Moth}: A pale, silent insect that clings to the walls of the upper corridors. It feeds on the lichen that grows between the stones. Not dangerous, but its wings reflect light in a way that can cause brief disorientation. Practical takeaway: If you see them clustering, turn around. They avoid the deeper chambers.
    \item \textbf{Pressure Eel}: A translucent, limbless creature that slides through cracks in the stone. It senses movement and heat. It does not attack, but its passage can trigger old mechanisms---floor tiles that tilt, doors that close. Practical takeaway: Watch the floor for ripples. If you see one, freeze. It will pass.
    \item \textbf{Warden Golem (Dormant)}: A guardian statue carved from obsidian, seated beside a sealed door. Its eyes are empty sockets. Tam touched its knee. It did not move. "It's not sleeping," he said. "It's waiting for the right question."
\end{itemize}

\section*{Journal Entry 4: The Thief and the Throne}
\textit{Day 3 (late). We found the central chamber. A domed space, larger than any hall I have seen, lit by a green light from no source. In the center, a sarcophagus of black stone. On its lid: the relic---a fist-sized eye, carved from jade, closed in dreaming. Beside it, a body.}

The relic hunter. His gear was scattered, his skull cracked against the base of the sarcophagus. He had fallen, or been pushed. His hand still reached for the stone.

Before we could approach, the light shifted. A figure stood beside the sarcophagus---translucent, robed, its face smooth and featureless. A ghost.

\subsection*{The Ghost's Account (Paraphrased from Tam's Later Summary)}
\textit{"The stone was mine. It was a gift from my mother, who was the keeper of this place before me. I was its guardian. Then treasure hunters came---not from Grey Hollow, but from the lowlands. They killed me and took the stone. I do not know what became of them. But the stone was stolen again, and again, and again.}

\textit{Generations later, the villagers of Grey Hollow found it in a dead man's pack. They did not know its history. They took it to their shrine and prayed to it. I did not stop them. It was... good to be prayed to. Good to feel remembered.}

\textit{But the stone is not a relic of luck. It is a reliquary. It holds my name. Without it, I cannot pass into whatever comes after. I need it back. So I whispered to the beggar who slept in the cave entrance. I whispered until he could not refuse. He stole it for me. He brought it here. He tried to return it to my sarcophagus, but the Warden... the Warden did not recognize him. It struck him down.}

\textit{I did not want him to die. But I need the stone."}

\subsection*{Tam's Rule No. 14: The Dead's Claim}
\textit{"A ghost that asks for its own relic is not a monster. It is a creditor. The debt is older than the village. Older than the stone. You cannot negotiate away what was stolen from the dead."}

Tam stood in silence for a long minute. Then he spoke.

"The villagers need the stone to survive. You need it to pass on. Neither of you is wrong. But the stone cannot be in two places at once."

The ghost tilted its head.

"I will take the stone to the village," Tam said. "I will let them keep it for one generation. One lifetime. Then I will return it to you."

"Why would I trust you?" the ghost asked.

"Because I am not a thief. I am a walker of thresholds. I will leave my compass here, with you. It was my father's. You can hold it until I return."

The ghost considered. Then it nodded. "One generation. No longer. If you are late, I will walk to the village myself."

\section*{Journal Entry 5: The Return}
\textit{Day 6. We left the ziggurat with the jade eye in Tam's pack. The walls did not shift. The stairs counted true---eighty, eighty, then the hall, then the warm air of the entrance. The ghost's light faded behind us.}

The village elder wept when she saw the stone. Tam placed it in her hands.

"It stays here for one generation," he said. "Then someone must return it to the deep place. Someone who is not afraid of the dark. Someone who can count to eighty and not panic."

"I will go," I heard myself say.

Tam looked at me. "You know the way?"

"I marked it. In my notebook. By the pressure eels and the warden's knee and the stair that felt longer on the way down."

He nodded. "Then you're not a scholar. You're a ranger."

The villagers built a new shrine. They offered us food---proper portions this time, thick stew, fresh bread. The children laughed. The fields, I noticed, had already begun to green.

\subsection*{Tam's Final Rule of the Journey}
\textit{"A relic is not a tool. It is a promise made by the dead to the living. Break the promise, and the dead will whisper. Keep it, and they will wait. Either way, you owe them. Choose what you can pay."}

\subsection*{Bestiary of the Ziggurat (Final)}
\begin{itemize}
    \item \textbf{The Warden (Dormant)}: A golem that guards the central chamber. It does not attack unless you approach the sarcophagus without the correct offering. The correct offering? The ghost's name, spoken in the old tongue. Tam learned it from the ghost before we left. He would not tell me. "You don't need it," he said. "You're not coming back alone.")
\end{itemize}

\vspace{2em}
\textit{End of the Ziggurat Journal. We left Grey Hollow at dawn. The villagers lined the road. No one spoke. They just watched, holding the green light of the jade eye in their minds. I will return in a generation. I will bring my own compass.}

\noindent Moira of Stonehaven\\
\textit{The road east, the lowlands still two days away. The ink is low. The stone is safe.}


\end{document}
